\begin{pastexam}[label={px:ncu-109-dsa-09}]{中央資工所 109 學年度 第 9 題}
（是非；答對 1 分，答錯倒扣 0.5 分）

An arbitrary graph with G(V; E), with $|E|=|V|-1$ edges is a binary tree.\\
(A)~True. \qquad (B)~False.
\end{pastexam}

\begin{pastsoln}
答案：(B) False。

錯在兩個地方，舉一個反例就夠，下面兩個都給。
\begin{itemize}
\item \textbf{邊數對不一定是樹}。 樹是連通、沒有環的無向圖，有 $|V| - 1$ 條邊；但反過來，有 $|V| - 1$
條邊的圖，要再加上「連通」或「無環」其中一個條件才是樹（資結書定理~21.3）。 反例：$4$ 個點，
邊是三角形 $1$--$2$、$2$--$3$、$3$--$1$，第 $4$ 個點孤立。 $|E| = 3 = |V| - 1$，但有環又不連通，不是樹。
\item \textbf{是樹也不一定是二元樹}。 二元樹要有根，每個節點最多兩個子節點，而且分左右（資結書
第~7.1~節）。 反例：星形圖，一個中心連到另外 $4$ 個點，$|V| = 5$、$|E| = 4$，是樹。 不管拿哪一點當
根，中心都至少有 $3$ 個子節點（當根時 $4$ 個，不當根時 $4$ 個鄰點裡一個是父節點、$3$ 個是子節點），
不可能是二元樹。 何況一般的圖沒有指定根和左右，本來就不是二元樹這種結構。
\end{itemize}
\end{pastsoln}
